\section{Pressure on the shell}
\label{sec:pressure}

\subsection{Fluent wall pressure}

Fluent gives the wall pressure as a pressure coefficient, which depends on the Mach number but hardly on
the altitude:
\begin{equation}
C_p(x) = \frac{p_w(x) - p_\infty}{q}.
\end{equation}
At the stagnation point $C_p = (p_0 - p_\infty)/q$: \val{A-cp0-M147} at burnout and \val{A-cp0-M095}
at $M = 0.95$ (Figure~\ref{fig:cp}). The pressure then drops quickly over the spherical tip and the first
centimetres, where A, with its fuller LV-Haack front, keeps about twice the pressure of B. Further aft
both noses are close to $p_\infty$, and below it towards the shoulder, where the flow accelerates onto
the cylinder: lowest $C_p$ \val{A-cpmin-M147} (A) and \val{B-cpmin-M147} (B) at burnout,
\val{A-cpmin-M095} and \val{B-cpmin-M095} at $M = 0.95$.

\begin{figure}[H]
\centering
\begin{subfigure}[t]{0.8\linewidth}
\centering
\includegraphics[width=\linewidth]{fluent_cp_AB_adiab}
\caption{Burnout, $M = \val{pt-M147-M}$.}
\end{subfigure}

\medskip
\begin{subfigure}[t]{0.8\linewidth}
\centering
\includegraphics[width=\linewidth]{fluent_cp_AB_M0.95}
\caption{Transonic, $M = 0.95$.}
\end{subfigure}
\caption{Wall pressure coefficient along \zone{wall\_nose}, A (red) and B (blue), Fluent XY plots. The
$M = 1.20$ run lies between the two and is in the data tables.}
\label{fig:cp}
\end{figure}

\begin{figure}[H]
\centering
\begin{subfigure}[t]{0.49\linewidth}
\centering
\includegraphics[width=\linewidth]{fluent_mach_A_adiab}
\caption{Nose of A.}
\end{subfigure}\hfill
\begin{subfigure}[t]{0.49\linewidth}
\centering
\includegraphics[width=\linewidth]{fluent_mach_tip_A_adiab}
\caption{Tip of A.}
\end{subfigure}
\caption{Mach number at burnout (Fluent). The bow shock stands just ahead of the spherical tip, with a
small subsonic region behind it where the stagnation pressure acts.}
\label{fig:mach}
\end{figure}

\subsection{From the Fluent runs to the whole flight}

At any time $t$ of the trajectory, $C_p$ is interpolated linearly in $M$ between the three runs, and
\begin{equation}
p_w(x, t) = p_\infty(t) + C_p\bigl(x, M(t)\bigr)\,q(t), \qquad
\Delta p(x, t) = p_w(x, t) - p_{in}(t),
\end{equation}
with $p_{in}$ of Equation~\eqref{eq:dp}. Outside the Fluent range $C_p$ is held: below $M = 0.95$ (low
$q$, so small loads) and above $M = \val{pt-M147-M}$, which only the hot day reaches ($M = \val{A-hot-LC1-M}$).
There the stagnation pressure comes out \val{A-hot-p0check}\,\% (A) and \val{B-hot-p0check}\,\% (B)
below the Rayleigh pitot value; on the nominal trajectory the difference is \val{A-nominal-p0check}\,\%
and \val{B-nominal-p0check}\,\%.

Between $x = 3$ and \SI{15}{mm} Fluent's $C_p$ has a small sawtooth where the \SI{0.2}{mm} tip cells
of the mesh grow to the \SI{1}{mm} wall cells. It is smoothed with a \SI{2}{mm} moving average in
the figures and tables below; the integrated forces are not affected.

\subsection{Net pressure on the sealed nose}

\begin{figure}[H]
\centering
\includegraphics[width=\linewidth]{fig_dp_x}
\caption{Net pressure along the nose at the three load cases, (a) A and (b) B. LC1 and LC2 with
$p_{in}$ at $T_{pad}$; LC3 with both bounds of $p_{in}$ (thick: $T_{pad}$, thin: heated air). $x$ on a
logarithmic scale above \SI{20}{mm}; dotted: shoulder.}
\label{fig:dpx}
\end{figure}

\begin{figure}[H]
\centering
\includegraphics[width=\linewidth]{fig_dp_t}
\caption{Net pressure along the flight at four stations, (a) A and (b) B. Thick: $p_{in}$ at
$T_{pad}$; thin: heated inside air.}
\label{fig:dpt}
\end{figure}

\begin{table}[H]
\centering
\caption{Extremes of the net pressure over the flight at the four stations [kPa], with the time [s]
in brackets. Inward: $p_{in}$ at $T_{pad}$; outward: heated inside air.}
\label{tab:stations}
\small
\begin{tabular}{@{}llcccc@{}}
\toprule
 & & A nominal & A hot day & B nominal & B hot day\\
\midrule
Tip & inward & \val{A-nominal-sttip-max} (\val{A-nominal-sttip-tmax}) & \val{A-hot-sttip-max} (\val{A-hot-sttip-tmax}) & \val{B-nominal-sttip-max} (\val{B-nominal-sttip-tmax}) & \val{B-hot-sttip-max} (\val{B-hot-sttip-tmax})\\
 & outward & \val{A-nominal-sttip-min} (\val{A-nominal-sttip-tmin}) & \val{A-hot-sttip-min} (\val{A-hot-sttip-tmin}) & \val{B-nominal-sttip-min} (\val{B-nominal-sttip-tmin}) & \val{B-hot-sttip-min} (\val{B-hot-sttip-tmin})\\
\SI{50}{mm} & inward & \val{A-nominal-st50-max} (\val{A-nominal-st50-tmax}) & \val{A-hot-st50-max} (\val{A-hot-st50-tmax}) & \val{B-nominal-st50-max} (\val{B-nominal-st50-tmax}) & \val{B-hot-st50-max} (\val{B-hot-st50-tmax})\\
 & outward & \val{A-nominal-st50-min} (\val{A-nominal-st50-tmin}) & \val{A-hot-st50-min} (\val{A-hot-st50-tmin}) & \val{B-nominal-st50-min} (\val{B-nominal-st50-tmin}) & \val{B-hot-st50-min} (\val{B-hot-st50-tmin})\\
Mid-length & outward & \val{A-nominal-stmid-min} (\val{A-nominal-stmid-tmin}) & \val{A-hot-stmid-min} (\val{A-hot-stmid-tmin}) & \val{B-nominal-stmid-min} (\val{B-nominal-stmid-tmin}) & \val{B-hot-stmid-min} (\val{B-hot-stmid-tmin})\\
Shoulder & outward & \val{A-nominal-stshoulder-min} (\val{A-nominal-stshoulder-tmin}) & \val{A-hot-stshoulder-min} (\val{A-hot-stshoulder-tmin}) & \val{B-nominal-stshoulder-min} (\val{B-nominal-stshoulder-tmin}) & \val{B-hot-stshoulder-min} (\val{B-hot-stshoulder-tmin})\\
\bottomrule
\end{tabular}
\end{table}

\begin{itemize}
\item \textbf{Inward} loads exist only on the front of the nose: the tip cap at burnout and, for A,
the first centimetres (\val{A-nominal-st50-max}\,kPa at \SI{50}{mm}, around $M = 1.2$). From mid-length
to the shoulder the net pressure is never inward.
\item \textbf{Outward} loads grow along the whole climb as $p_\infty$ falls, and reach about
\SI{70}{kPa} on the whole shell at apogee. Heating of the inside air adds a few kPa during the coast
(thin lines in Figure~\ref{fig:dpt}); at the shoulder of B this gives the largest outward value,
\val{B-hot-stshoulder-min}\,kPa at $t = \val{B-hot-stshoulder-tmin}$\,s on the hot day.
\item Over the shell the two designs differ only at the front; at apogee they are identical, because
there $\Delta p = p_\infty - p_{in}$ does not depend on the shape.
\end{itemize}
The full distributions $p_w(x)$ and $\Delta p(x)$ at LC1--LC3, and $\Delta p(t)$ at the four stations,
are in the data tables (Appendix~\ref{app:data}).
